Introduction
A central goal of mechanistic interpretability is to decompose a neural network's parameters into a small number of meaningful units, atoms, each of which corresponds to something the model is doing. The hope is that a network's computation, however dense and distributed in its raw form, can be expressed as a sum of simpler mechanisms that admit human-readable descriptions. The viability of this programme depends both on whether such atoms exist in a recoverable form, and on whether the methods used to extract them recover the right ones.
A recent line of work shifts the object of decomposition from activations to parameters, decomposing a network's weights into a sum of simpler components, each taken to correspond to a concrete atom of computation. Methods guided by the principle of minimum description length aim to recover components that are faithful to the original weights, minimal in the number required per input, and individually as simple as possible. Stochastic Parameter Decomposition (SPD) [Bus+25] is one such method, decomposing networks into rank-one subcomponents and learning a per-component importance function that estimates each atom's causal contribution to a given input. The present work is a direct extension of SPD.
Whether such a decomposition recovers the computational primitives the model actually uses depends on what the ground-truth decomposition is taken to be. For models with independent features the answer is comparatively unambiguous: one atom per feature. The question becomes substantially harder when the features themselves are dependent, a regime that is empirically attested in real models. Engels et al. [Eng+24] develop a rigorous notion of irreducible multi-dimensional features and show that language models represent inherently cyclical concepts, including days of the week and months of the year, as circles in activation space whose individual coordinates are not independent features but components of a single object the model manipulates as a unit.
Networks of this kind admit multiple decompositions that are equally faithful to the weights but differ in the granularity of the units they expose. A coordinate-level decomposition is atomic in the strictest sense, but the structure that makes the coordinates meaningful as a group is left implicit, recoverable only by inspecting which atoms co-activate. A structure-level decomposition exposes that structure directly, at the cost of treating as a single unit an object that is internally composite, and risks obscuring the computation taking place inside that unit. A decomposition coarser than the model's computation conceals what happens inside its atoms; a decomposition finer than the model's structural relations conceals those relations between its atoms. Each direction loses something the other preserves. The Open Problems in Mechanistic Interpretability programme [Sha+25] identifies the question of how best to carve a network at its joints as a central open problem of the field; the present work takes up one specific instance of that question, in the setting of parameter decomposition.
Applied to dependent features, SPD does not, by default, produce the coordinate-level decomposition its rank-one formulation might suggest. Instead, it recovers a structure-level decomposition in which groups of dependent features share overlapping subcomponents — a clustering that the original SPD paper observes must otherwise be performed post hoc [Bus+25], and which subsequent work extending SPD to transformers [Chr+25] continues to treat as a separate downstream step. The present work shows that the finer coordinate-level decomposition is reachable on the same model when softmax-and-attention competition is introduced over subcomponents, forcing each atom to justify itself against the alternatives and sharpening the allocation toward one atom per feature. The position taken here is that the choice between the two readings is a substantive methodological commitment rather than a problem the data can settle, and that operationalising the choice within SPD requires developing both readings as deliberate proposals. Doing so makes the structural commitment a first-class object of the decomposition rather than something inferred from it, and forces the practitioner to make deliberately what is otherwise made implicitly through the form of the regulariser.